\section{Introduction}
\label{sec:intro}

Watermarking is one of the few mechanisms that lets a model provider later recognise text its model produced without storing every output. The most widely studied family, introduced by \citet{kirchenbauer2023watermark} and commonly called KGW, hashes a secret key together with the preceding token to pseudo-randomly select a ``green'' fraction $\gamma$ of the vocabulary, adds a bias $\delta$ to the logits of green tokens, and detects the watermark with a one-proportion $z$-test on the count of green tokens. Variants of this idea now ship in widely used open-source libraries \citep{wolf2020transformers,pan2024markllm} and are evaluated in dedicated benchmarks \citep{tu2024waterbench,piet2023markmywords}.

The key is treated as a free parameter. Because the partition is produced by a hash, any valid key is assumed to be as good as any other, and quality evaluations of watermarking typically report results for a single, fixed key. This assumption is plausible for prose, where many surface forms are acceptable. It is less obviously safe for code. A program's security can depend on a single token, such as \texttt{snprintf} versus \texttt{sprintf} or \texttt{yaml.safe\_load} versus \texttt{yaml.load}, and LLM-generated code is already frequently insecure without any watermark \citep{pearce2022asleep,bhatt2023cyberseceval,peng2025cweval}. For a fixed key and context, KGW applies the same deterministic preference every time that context recurs. If a key happens to place insecure idioms on the green side of recurring contexts, its bias could recur in many generated programs. Whether this happens is an empirical question that, to our knowledge, has not been tested for code security.

\paragraph{Two questions, not one.} Answering it requires separating two estimands that are easily conflated. The \emph{average effect} of KGW compares watermarked and matched unwatermarked generation and measures the price of watermarking. \emph{Between-key heterogeneity} compares valid keys with each other and tests the interchangeability assumption itself. A large spread between the best and worst of a few keys is not evidence of heterogeneity unless it exceeds what sampling noise and task clustering produce, and a heterogeneity claim is only meaningful if the watermark manipulation was actually applied and the security oracle can discriminate.

\paragraph{This paper.} We study a single open model, Qwen2.5-Coder-7B-Instruct \citep{hui2024qwen25coder}, with the KGW implementation in Hugging Face Transformers. Our evidence is preliminary and we report it as such. The pilot shows large key-associated variation in \emph{functional} correctness. It does not show key-dependent \emph{security}, and its detector diagnostics show that the watermark manipulation is not yet validated for short programs. We think these diagnostic and negative findings are useful in their own right: they show which validity threats a key-audit study has to rule out before any key can be called harmful.

\paragraph{Contributions.}
\begin{itemize}
\item \textbf{Formulation.} We define the average KGW effect and between-key heterogeneity as separate estimands over a paired task$\times$seed design. We prove that between-key contrasts are invariant to the shared unwatermarked baseline (Proposition~\ref{prop:baseline}) and derive a first-order minimum output length for per-output KGW detection (Proposition~\ref{prop:length}).
\item \textbf{Protocol and artifact.} We describe a fail-closed audit pipeline: deterministic key derivation, provenance-stamped resumable generation, a networkless sandbox for executing generated code, deterministic C/Python oracles over 12 CWE families, and preregistered manipulation and go/no-go gates.
\item \textbf{Pilot evidence.} A 7,200-generation calibration shows a monotone bias dose response. A 792-program ten-key pilot shows functional rates from 12.5\% to 70.8\% across keys, near-constant security labels, and weak per-output detection (21/720 outputs at $z\!\ge\!4$; five of ten keys with negative mean $z$).
\item \textbf{Validity analysis.} We list the threats that currently prevent a security claim and the concrete gates that address them, so that the planned 100-key study can produce a falsifiable result.
\end{itemize}
